\documentclass[10pt,a4paper]{amsart}
\usepackage[margin=19mm]{geometry}
\usepackage{amsmath,amssymb,mathtools}
\usepackage[scr=rsfs,cal=euler]{mathalfa}
\usepackage{xfrac}
\usepackage[unicode,hidelinks,bookmarks=false]{hyperref}
\newcommand{\ie}{i.e.\ }
\newcommand{\cf}{cf.\ }
\newcommand{\resp}{respectively\ }
\newcommand{\Subsection}{\subsection}
\providecommand{\nobreakdash}{\nobreak\hbox{-}}
\setlength{\emergencystretch}{2em}
\multlinegap0pt
\usepackage{bm}
\usepackage{bbm}
\usepackage{dsfont}
\usepackage{xspace}
\usepackage{cancel}
\usepackage{mathtools}
\usepackage{amsthm}
\makeatletter
\@ifundefined{thm}{\newtheorem{thm}{Thm}}{}
\@ifundefined{prop}{\newtheorem{prop}[thm]{Proposition}}{}
\makeatother
\makeatletter
\expandafter\ifx\csname emphit\endcsname\relax
\newenvironment{emphit}{\begin{itemize}\em}{\end{itemize}}
\fi
\expandafter\ifx\csname prf\endcsname\relax
\newenvironment{prf} {{\bf Proof.}}{\hfill $\Box$}
\fi
\makeatother
\title[Frame properties induced by iteration of a multiplication op]{Frame properties induced by iteration of a multiplication operator on Hardy spaces}
\author{Jahangir Cheshmavar}
\date{Source: arXiv:2509.04879v1 (CC BY 4.0)}
\begin{document}
\maketitle

\noindent\emph{Source and licence.} Frame properties induced by iteration of a multiplication operator on Hardy spaces. Version arXiv:2509.04879v1 (CC BY 4.0), distributed under the Creative Commons Attribution 4.0 International licence, as recorded in the arXiv licence metadata for this exact version and shown on the abstract page https://arxiv.org/abs/2509.04879v1. Full rights notice: licenses/arxiv-2509.04879-CC-BY-4.0.txt.\par\smallskip

\noindent\textit{Editorial scope.} Counted unit: the Prop whose source label is \texttt{Prop.1} in the article (source0.tex lines 263--269, source label \texttt{Prop.1}), delivered together with the article's own complete original proof of it (source0.tex lines 270--366, one \texttt{proof} environment of 97 lines). No mathematics is written, repaired or re-derived: statement and proof are the article's own text line for line. Required context retained uncounted: none. Every definition, display and label of those lines is retained unchanged. Omitted from this package and not redistributed: 1--262, 367--948. Nothing inside a retained excerpt is omitted or altered: every retained body is delivered exactly as the article's lines are written, so no omission receipt of any kind accompanies this package. Citations: the retained excerpts carry no citation command at all, so no bibliography entry of the article is transcribed and no reference is redistributed. Reference adaptation: none was needed and none was made; every reference inside the delivered unit resolves inside it, so no unresolved reference or citation is produced. Printed slips kept literal: no printed word, symbol, hypothesis or number of the counted statement or of its proof was altered. Layout and class: the article's own class is \texttt{amsart} with packages including amsopn, amssymb, amsaddr, graphicx, xy, amscd, color, hyperref, amsmath, amsthm, amsfonts, verbatim, latexsym; the wrapper is the lane's pinned \texttt{amsart} editorial wrapper at 10pt on A4, which restates the theorem-like environments the retained text needs and adds the article's own macro definitions that the retained text needs (none), with extra wrapper packages (bm, bbm, dsfont, xspace, cancel, mathtools). The delivered numbering is the unit's own. Assets: no retained span contains a figure environment, a graphics-inclusion command, a PDF-page inclusion, a table or a TikZ picture; no image file is copied into this package, the package declares no asset dependency and no image file is redistributed with it. Licence: version arXiv:2509.04879v1 of this article is distributed under the Creative Commons Attribution 4.0 International licence, as recorded in the arXiv licence metadata for this exact version and shown on the abstract page https://arxiv.org/abs/2509.04879v1. A full rights notice is delivered at licenses/arxiv-2509.04879-CC-BY-4.0.txt. Characters: every retained excerpt is pure ASCII, so the delivered engine, XeLaTeX with its default Latin Modern text font, reports no missing character.
\begin{prop} \label{Prop.1}
For any $\phi \in H^{\infty}(\mathbb{D})$, let $T_{\phi}:
H^2(\mathbb{D})\rightarrow H^2(\mathbb{D})$ be the associated
multiplication operator. If the system
$\{T_{\phi}^nf\}_{n=0}^{\infty}=\{ \phi^n f \}_{n=0}^\infty $ is a
frame for $H^2(\mathbb{D})$, then $\phi$ is an inner function.
\end{prop}

\begin{proof}
Suppose $\phi \in H^\infty(\mathbb{D})$ is not inner, meaning
$|\phi(e^{i\theta})|\neq 1$ on a set of positive measure.
Typically, for non-inner functions, $|\phi(e^{i\theta})|<1$ (since
$\|\phi\|_{H^{\infty}}\leq 1$ often in such contexts , or we
normalized), so we have $|\phi(e^{i\theta})|<1 \,\ \text{on a set
} E
\subset \mathbb{T} \text{ of positive measure}$.\\
Claim: $\|\phi ^n f\| \rightarrow 0$, as $n \rightarrow \infty$
for all $f \in H^2(\mathbb{D})$.

For the proof of claim, we not that the set $E$ satisfies $|E| >
0$ (i.e., $\phi$ is not unimodular a.e. on $\mathbb{T}$). Then for
all $n \geq 1$, $|\phi(e^{i\theta})|^n \to 0 \quad \text{uniformly
on } E$. Now let $f \in H^2(\mathbb{D})$, $f \neq 0$. Then the
boundary function $f(e^{i\theta})$ belongs to $L^2(\mathbb{T})$.
So
\begin{eqnarray*}
\|\phi^n f\|_{H^2}^2 = \int_0^{2\pi} |\phi(e^{i\theta})|^{2n}
|f(e^{i\theta})|^2 \, \frac{d\theta}{2\pi}.
\end{eqnarray*}
Now we split the integral over two parts:
\begin{itemize}
\item[(i)] On $E$, $|\phi(e^{i\theta})| < 1$, and the integrand
decays exponentially in $n$. \item[(ii)] On $\mathbb{T} \setminus
E$, $|\phi(e^{i\theta})| \leq 1$, so the integrand is bounded.
\end{itemize}
  Hence,
\begin{eqnarray*}
\|\phi^n f\|_{H^2}^2 = \int_E |\phi(e^{i\theta})|^{2n}
|f(e^{i\theta})|^2 \, \frac{d\theta}{2\pi} + \int_{\mathbb{T}
\setminus E} |\phi(e^{i\theta})|^{2n} |f(e^{i\theta})|^2 \,
\frac{d\theta}{2\pi}
\end{eqnarray*}
Let's denote:
\begin{eqnarray*}
I_n^{(1)} := \int_E |\phi(e^{i\theta})|^{2n} |f(e^{i\theta})|^2 \,
\frac{d\theta}{2\pi}, \quad I_n^{(2)} := \int_{\mathbb{T}
\setminus E} |\phi(e^{i\theta})|^{2n} |f(e^{i\theta})|^2 \,
\frac{d\theta}{2\pi},
\end{eqnarray*}
Then Since $|\phi(e^{i\theta})| < 1$ on $E$, there exists $r < 1$
such that $|\phi(e^{i\theta})| \leq r$ for a.e. $\theta \in E$.
Thus,
\begin{eqnarray*}
I_n^{(1)} \leq r^{2n} \|f\|_{L^2(E)}^2 \to 0 \quad \text{as } n
\to \infty.
\end{eqnarray*}
On $\mathbb{T} \setminus E$, $|\phi(e^{i\theta})| \leq 1$, so
$I_n^{(2)} \leq \|f\|_{L^2(\mathbb{T})}^2$, but can also go to
zero if $|\phi(e^{i\theta})| < 1$ on a set of positive measure.
But in any case,
\begin{eqnarray*}
\liminf_{n \to \infty} \|\phi^n f\|_{H^2}^2 \leq \lim_{n \to
\infty} I_n^{(1)} = 0.
\end{eqnarray*}
Therefore, $\|\phi^n f\|_{H^2} \to 0 \quad \text{as } n \to
\infty$. Of course, by dominated convergence theorem, also we
conclude
\begin{eqnarray*}
\lim_{n\rightarrow\infty}\|\phi^n f\|_{H^2}^2 = \int_0^{2\pi}
\lim_{n\rightarrow\infty}|\phi(e^{i\theta})|^{2n}
|f(e^{i\theta})|^2 \, \frac{d\theta}{2\pi}=0.
\end{eqnarray*}
Therefore, the vectors $\phi^n f$ are becoming arbitrarily small
in norm, and their span cannot remain "rich" enough to approximate
all of $H^2$, and hence cannot form a frame. For example, take $g
= f$. Then
\begin{eqnarray*}
\sum_{n=0}^\infty |\langle f, \phi^n f \rangle|^2 \leq \sum \|f\|
\|\phi^n f\|^2 = \|f\|^2 \sum \|\phi^n f\|^2 \to 0,
\end{eqnarray*}
since $\|\phi^n f\| \to 0$. This contradicts the existence of a
positive lower bound $A \|f\|^2$. More precisely, suppose
$\{\phi^n f\}_{n=0}^\infty$ is a frame for $H^2(\mathbb{D})$.
Then, in particular, there exists a lower frame bound $A>0$ such
that for all $g \in H^2(\mathbb{D})$,
\begin{eqnarray*}
\sum_{n=0}^\infty |\langle g, \phi^n f \rangle|^2 \geq A \|g\|^2.
\end{eqnarray*}
Let's test this inequality using some carefully chosen $g$. Let $g
= \phi^N f$ for some large $N$. Then
\begin{eqnarray*} \langle g, \phi^n f \rangle = \langle \phi^N f,
\phi^n f \rangle = \langle f, \phi^{n - N} f \rangle.
\end{eqnarray*}
 More concretely:
\begin{eqnarray*}
\sum_{n=0}^\infty |\langle \phi^N f, \phi^n f \rangle|^2 =
\sum_{n=0}^\infty |\langle f, \phi^{n - N} f \rangle|^2.
\end{eqnarray*}
This sum is dominated by $\|f\|^4$ and decays as $N \to \infty$,
because $f$ becomes orthogonal to $\phi^n f$ for large $n$, due to
decay of $\phi^n f$. Also, since $\|\phi^N f\| \to 0$, the
left-hand side becomes arbitrarily small, while the right-hand
side requires $A \|\phi^N f\|^2 \not\to 0$. Contradiction and
hence, the frame inequality fails if $\phi$ is not inner.
\end{proof}

\end{document}
